%------------------------------------------------------------------------------%
\section{Indução Finita}
\label{sec:Inducao-Finita}

Demonstração por \emph{Indução Finita}\index{Indução!finita}
ou
\emph{Indução Matemática}\index{Indução!matemática}
é uma técnica para demonstrar afirmações que dependem de um índice que segue para
o infinito, por exemplo, $n\in\N$.

Para mostrar que uma afirmação, $P_n$, é verdadeira para todo $n$ devemos provar que:
\begin{enumerate}
	\item ela é verdadeira para o primeiro $n$, por exemplo, se o índice começa em $n=1$
    temos que mostrar que $P_1$ é verdade;
	\item se ela for verdadeira para $n$ então também será verdadeira para o próximo
    índice, $n+1$, isso é
    \[
      P_n \Rightarrow P_{n+1} \,.
    \]
\end{enumerate}

O exemplo a seguir ilustra a técnica provando a
\emph{Desigualdade de Bernoulli}%
\index{Desigualdade!Bernoulli}%
\index{Bernoulli!desigualdade}.

\begin{example}
  \label{ex:Desigualdade_de_Bernoulli}
  Demonstrar a Desigualdade de Bernoulli
  \[
    (1+x)^n \geq 1 + n x
  \]
  para todo $x>-1$ e $n$ inteiro não negativo.

  \solution
  Primeiro vamos identificar qual é o índice envolvido e qual a afirmação
  que deve ser verificada para cada valor do índice
  \begin{gather*}
    n = 0, 1, 2, 3, \dots \\
    P_n\colon (1+x)^n \geq 1 + n x
  \end{gather*}
  Provamos agora o primeiro caso, $n=0$, por simples substituição
  \[
    (1+x)^0 = 1 \geq 1 + 0x = 1
  \]
  O próximo passo é assumir que a afirmação é verdadeira para $P_n$,
  que chamamos de \emph{Hipótese de Indução}\index{Hipótese de indução},
  e provar que ela é verdadeira para $P_{n+1}$.
  Neste exemplo queremos provar que
  \[
    (1+x)^n \geq 1 + nx \quad\Rightarrow\quad (1+x)^{n+1} \geq 1 + (n+1) x
  \]
  Começamos manipulando o lado esquerdo da desigualdade que queremos demonstrar
  \[
    (1+x)^{(n+1)} = (1+x)^{n}(1+x)
  \]
  Como $(1+x)>0$, pois por hipótese $x>-1$, podemos usar a hipótese de indução,
  $(1+x)^n \geq 1 + nx$, para escrever
  \[
    (1+x)^{(n+1)}
    \geq (1 + nx)(1 + x)
    = 1 + x + nx + nx^2
  \]
  Rearranjando temos
  \[
    (1+x)^{(n+1)}
    \geq 1 + (1 + n)x + nx^2
  \]
  Como $nx^2$ é sempre não negativo, podemos remove-lo mantendo a relação de desigualdade
  \[
    (1+x)^{(n+1)}
    \geq 1 + (1 + n)x
  \]
  Provamos assim que $P_n\Rightarrow P_{n+1}$, o que completa a prova por indução finita
  e garante que a desigualdade é verdadeira para todo $n\geq 0$.
\end{example}






%------------------------------------------------------------------------------%
